\documentclass[10pt,a4paper]{amsart}
\usepackage[margin=19mm]{geometry}
\usepackage{amsmath,amssymb,mathtools}
\usepackage[scr=rsfs,cal=euler]{mathalfa}
\usepackage{xfrac}
\usepackage[unicode,hidelinks,bookmarks=false]{hyperref}
\newcommand{\ie}{i.e.\ }
\newcommand{\cf}{cf.\ }
\newcommand{\resp}{respectively\ }
\newcommand{\Subsection}{\subsection}
\providecommand{\nobreakdash}{\nobreak\hbox{-}}
\setlength{\emergencystretch}{2em}
\multlinegap0pt
\usepackage{mathtools}
\usepackage{tcolorbox}
\usepackage{float}
\usepackage{algorithm}\usepackage{algpseudocode}
\usepackage{bm}
\usepackage[shortlabels]{enumitem}
\newtheorem{lemma}{Lemma}
\title{Collective Annealing by Switching Temperatures:\\ a Boltzmann-type framework}
\author{Fr\'ed\'eric Blondeel, Lorenzo Pareschi, Giovanni Samaey}
\date{Source: arXiv:2512.13522v2 (CC BY 4.0)}
\begin{document}
\maketitle

\noindent\emph{Source and licence.} Collective Annealing by Switching Temperatures:\\ a Boltzmann-type framework. Version arXiv:2512.13522v2 (CC BY 4.0), distributed by arXiv under the Creative Commons Attribution 4.0 International licence, http://creativecommons.org/licenses/by/4.0/, as recorded in the arXiv licence metadata for this exact version and shown on the abstract page https://arxiv.org/abs/2512.13522v2. Full licence notice: \texttt{licenses/arxiv-2512.13522-CC-BY-4.0.txt}.\par\smallskip

\noindent\textit{Editorial scope.} Counted unit: the article's lemma lemma:moment2 (source0.tex lines 707--716). The article's complete original proof of it is delivered with it (source0.tex lines 718--756, one \texttt{proof} environment of 39 lines). No mathematics is written, repaired or re-derived: the statement and the proof are the article's own lines, in the article's own notation, and no retained line is substituted or re-flowed.  No further source-local context is needed: every cross-reference of every retained line resolves inside the two delivered spans.  Nothing was omitted from any retained span, so the package carries no omission record.  No retained line contains a citation, so the wrapper carries no bibliography and no bibliography entry.  No retained line contains a figure environment, an image-inclusion command or drawing code, so no image is delivered and the package declares no asset dependency.  Editorial formatting only, in addition: an AMS \texttt{amsart} preamble that restates the article's own declarations and macros used by the retained lines, namely the environments \texttt{lemma} on separate counters, together with the standard packages those lines need.  The retained environments print in this document as Lemma 1 in order of appearance, whereas the article prints the same items with the numbering of its own full document.  The complete native source of the article as distributed in the arXiv e-print for this exact version is bundled unchanged as \texttt{sources/191034/source0.tex}.
\begin{lemma}[Evolution of the second moment]
\label{lemma:moment2}
The evolution of $
M_2:=\int_{\Omega}fT^2d\omega$
satisfies
 \begin{equation*}
\frac{dM_2}{dt}=\int_{\Omega\times\Omega}\left[(\mu-\lambda)(T^2-T_*^2)+(\lambda^2+\mu^2-\lambda-\mu)(T-T_*)^2\right]\chi_\mathcal{F}(x,x_*,T,T_*)ff_*d\omega d\omega_*.
\end{equation*}
for $\lambda,\mu\in [0,1]$, $\mu\leq\lambda$ and $\sigma(x,x_*,T,T_*)=0$. 
\end{lemma}

\begin{proof}
We start from the post–interaction temperature update
\[
T' = T - \lambda (T - T_*)\chi_\mathcal{F}(x,x_*,T,T_*) - \mu (T - T_*)\chi_\mathcal{F}(x_*,x,T_*,T),
\]
and compute
\[
\begin{split}
\mathbb{E}[T'^2 - T^2]
= &-2\lambda T(T-T_*)\chi_\mathcal{F}(x,x_*,T,T_*) 
  -2\mu T(T-T_*)\chi_\mathcal{F}(x_*,x,T_*,T)\\
  &+ \lambda^2 (T-T_*)^2\chi_\mathcal{F}(x,x_*,T,T_*) 
  + \mu^2 (T-T_*)^2\chi_\mathcal{F}(x_*,x,T_*,T),
  \end{split}
\]
where we used  
\[
\chi_\mathcal{F}(x,x_*,T,T_*)\chi_\mathcal{F}(x_*,x,T_*,T)=0,\qquad \chi^2_\mathcal{F}(x,x_*,T,T_*)=\chi_\mathcal{F}(x,x_*,T,T_*).
\]
Integrating against $ff_*$ on $\Omega\times\Omega$ yields
\begin{equation*}
\begin{split}
\frac{dM_2}{dt}&=\int_{\Omega\times\Omega}\left[-2\lambda T(T-T_*)-2\mu T_*(T_*-T)\right]\chi_\mathcal{F}(x,x_*,T,T_*)ff_*d\omega d\omega_*\\
&+\int_{\Omega\times\Omega}(\lambda^2+\mu^2)(T-T_*)^2\chi_\mathcal{F}(x,x_*,T,T_*)ff_*d\omega d\omega_*.
\end{split}
\end{equation*}
Expanding the first integral and rearranging yields
 \begin{equation*}
\begin{split}
\frac{dM_2}{dt}&=\int_{\Omega\times\Omega}(\mu-\lambda)(T^2-T_*^2)\chi_\mathcal{F}(x,x_*,T,T_*)ff_*d\omega d\omega_*\\
&+\int_{\Omega\times\Omega}(\lambda^2+\mu^2-\lambda-\mu)(T-T_*)^2\chi_\mathcal{F}(x,x_*,T,T_*)ff_*d\omega d\omega_*.
\end{split}
\end{equation*}

The first integral is either negative or zero if $\mu\leq\lambda$. Since $\lambda(\lambda-1)+\mu(\mu-1)\leq 0$ for $\lambda,\mu\in[0,1]$, the second integral is non-positive, and it is strictly negative if the support of
$\chi_\mathcal{F}(x,x_*,T,T_*)$ has positive measure and $\lambda,\mu\in(0,1)$.
Thus, in the absence of interaction noise, \(M_2\) is non-increasing under the
stated assumptions. 
\end{proof}

\end{document}
